\documentclass{sprawozdanie-agh}

\usepackage[utf8]{inputenc}
\usepackage{matlab-prettifier}
\usepackage{siunitx}
\usepackage{booktabs,caption,ragged2e}

\makeatletter

% ========= Początek Dokumentu ========= 

\begin{document}


% ========= Strona tytułowa ========= 
\przedmiot{Teoria Sterowania II: Ćwiczenia Laboratoryjne}
\tytul{Wytyczne dla Ćwiczenia 1 }
\podtytul{Portrety Fazowe Systemów Liniowych}
\kierunek{Automatyke i Robotyka}
\autor{Dariusz Cieślar, Maciej Różewicz}
\data{Kraków, 3 marca 2026}

\stronatytulowa{}

% ========= Treść Właściwa =========

\section{Przebieg ćwiczenia}
\subsection{Zadanie}

\par
W świecie STATESPACE grasuje dziewięcioro przestępców. Ich działania nie pozostają bez śladu. Spędziliśmy cały zeszły semestr, żeby zebrać na nich dowody. Niestety w archiwum mamy niezły bałagan: wszystko się pomieszało i teraz musimy to uporządkować.
\vspace{0.5cm}
\par
Podzielimy się na ekipy. Każda z ekip otrzyma inny zestaw czterech dowodów. 
\vspace{0.5cm}
\par
Zadaniem każdej z ekip jest odnalezienie wszystkich dowodów (wśród   dostępnych w całej grupie) dotyczących macierzy, dla których ekipa ta otrzymała przynajmniej jeden dowód.
\vspace{0.5cm}
\par
\textbf{Na przykład:}
Ekipa 3 otrzymała dowody 9, 10, 11 i 12. W trakcie zajęć ekipa znalazła, że dowód 9 opisuje tą samą macierz co dowód 19 i 22. Dowód 10 opisuje tą samą macierz co już posiadany dowód 12, ale u innych grup znaleziono też pasujące dowody 4 i 27. Wreszcie dowód 11 opisuje tą samą macierz co odnalezione u innych ekip w grupie dowody 3, 5 i 18.
\vspace{0.5cm}
\par
Zatem poprawne rozwiązanie ćwiczenia dla ekipy 3 to:
\begin{table}[h]
\caption{Rozwiązanie zadania}
\centering
\begin{tabular}{| m{2cm} | m{3cm}| m{3cm} | m{3cm} |}
\hline
Zestaw & Otrzymany Dowód & Pasujące Dowody & Postać Macierzy  \\
\hline
A      & 9               & 19, 22          & {[}1 0; 2 -1{]}  \\
\hline
B      & 10, 12          & 4, 27           & {[}2 0; 0 -2{]}  \\
\hline
C      & 11              & 3, 5, 18        & {[}0 1; -3 -8{]} \\
\hline
\end{tabular}
\label{tab:przyklad}
\end{table}

\subsection{Przykłady i umiejętności do zdobycia}
W folderze \textit{przykłady/01_portrety_fazowe} znajdziesz skrypty .mlx, które możesz uruchomić w MATLABie. Zapoznaj się z ich zawartością - będą one pomocne w wykonaniu ćwiczenia. 

Zwróć uwagę jak przebiegi czasowe dla wybranych warunków początkowych mają się do trajektorii fazowej dla tych warunków początkowych.


W szczególności upewnij się, że potrafisz:
\begin{description}
\item[1] Przeprowadzić symulacje odpowiedzi swobodnej dla układów LTI opisanych macierzami stanu 2x2.
\item[2] Wykreślić uzyskane w rezultacie symulacji trajektorie fazowe na płaszczyźnie stanu.
\item[3] Wykreślić na płaszczyźnie zespolonej pozycje wartości własnych .
\item[4] Wykreślić potret fazowy dla zadanej macierzy o rozmiarze 2x2.
\item[5] Uzupełnić portet fazowy o:
\begin{itemize}
\item wykreślenie podprzestrzeni niezmienniczych,
\item wykreślenie rzeczywistych wektorów własnych (o ile istnieją),
\item zaznaczenie punktów równowagi.
\end{itemize}

\item[6] Zmodyfikować daną macierz stanu przez przekształcenie po postaci Jordana oraz przekształcenie przez podobieństwo, np. do postaci Frobeniusa. Zweryfikować jak te przekształcenia wpływają na postać portretu.

\item[7] Wykreślić przebiegi czasowe zmiennych stanu na wykresie 2D i 3D

\item[8] Sprawdzić jak zamiana znaku wartości własnych (np. ujemne na dodatnie) wpływa na portret fazowy i przebiegi czasowe w odpowiedzi swobodnej.

\end{description}

\vspace{0.5cm}
\par
%\textbf{Zadanie dodatkowe}

%Zaproponuj fizyczną realizację układu drugiego rzędu tak, aby modyfikując wartości elemetów tego układu, można było uzyskać co najmniej 5 typów portretów fazowych.

%Wykreśl schemat oraz sformułuj model fizyczny tego układu. Wyznacz przykładowe wartości parametrów tego modelu dla każdego z 5 typów portetów fazowych.

\section{Wytyczne do sprawozdania}

Przygotuj sprawozdanie z ćwiczenia opisujące cel, przebieg i obserwacje.

W części opisującej przebieg ćwiczenia proszę koniecznie zawrzeć:    
\begin{enumerate}
        \item Tabelę z dopasowaniem do dowodów otrzymanych przez Twoją ekipę wszystkich pasujących dowodów w grupie
        \item Opis (każdego dopasowania z osobna) dlaczego Twoim zdaniem te znalezione u innych dowody pasują do dowodów otrzymanych przez Twoją ekipę. 
        \item Dla każdego wiersza w Twojej tabeli: portrety fazowe z naniesionymi trajektoriami pochodzącymi z dowodów w postaci przebiegów czasowych 
\end{enumerate}


%\textbf{Zadanie dodatkowe}
%Proszę przedstawić rozwiązanie zadania dodatkowego w postaci:
%\begin{itemize}
%    \item Schematu proponowanego układu
%    \item Równań różniczkowych opisujących model fizyczny tego układu 
%    \item Tabelę pokazującą parametry elementów układu, postać macierzy stanu, wartości własne dla tych parametrów oraz portety fazowe 
%\end{itemize}

\end{document}